% Generated draft: requires mathematical and source review.
\documentclass[11pt]{article}
\usepackage[margin=25mm]{geometry}
\usepackage{fontspec}
\setmainfont{Latin Modern Roman}
\usepackage{amsmath,amssymb,mathtools,unicode-math}
\setmathfont{Latin Modern Math}
\usepackage{enumitem,needspace,xurl,hyperref}
\setlength{\parindent}{0pt}
\setlength{\parskip}{5pt}
\setlength{\emergencystretch}{3em}
\newcommand{\N}{\mathbb{N}}
\newcommand{\Z}{\mathbb{Z}}
\newcommand{\Q}{\mathbb{Q}}
\newcommand{\R}{\mathbb{R}}
\newcommand{\C}{\mathbb{C}}
\newcommand{\F}{\mathbb{F}}
\newcommand{\A}{\mathbb{A}}
\newcommand{\PP}{\mathbb{P}}
\newcommand{\E}{\mathbb{E}}
\newcommand{\eps}{\varepsilon}
\newcommand{\dd}{\,\mathrm{d}}
\begin{document}
\section*{UnsolvedMath ID: 700616}
\textbf{Source problem code:} catalogue entry
\subsection{Definitions and mathematical statement}
Let \(p\) be a prime. A \(p\)-subgroup of a group \(G\) is a subgroup whose elements all have finite order equal to a power of \(p\). A Sylow \(p\)-subgroup means a \(p\)-subgroup maximal under inclusion. Write \(\operatorname{Syl}_p(G)\) for the set of these subgroups. A group is periodic if every element has finite order.

The source asks: suppose that \(G\) is a ``(periodic) \(p\)-conjugacy biprimitively finite group'' in the sense of Kourovka Notebook Problem 6.59, and suppose that
\[
\exists S\in\operatorname{Syl}_p(G)\quad |S|<\infty.
\]
Must it follow that
\[
\forall P,Q\in\operatorname{Syl}_p(G)\quad
\exists g\in G\quad g^{-1}Pg=Q?
\]
The defining condition for ``\(p\)-conjugacy biprimitively finite'' is not included in the supplied record, so that hypothesis remains an unresolved reference to Problem 6.59. The parentheses around ``periodic'' are also retained from the source: the supplied record does not determine whether periodicity is required or whether periodic and unrestricted versions are both intended. Thus this is a conditional formulation pending clarification of the source hypothesis, rather than a complete standalone statement.
\subsection{Short English statement}
V. P. Shunkov asks whether a group satisfying the source's ``(periodic) \(p\)-conjugacy biprimitively finite'' hypothesis and having one finite Sylow \(p\)-subgroup must have all its Sylow \(p\)-subgroups conjugate. The technical hypothesis and the role of periodicity require clarification from the referenced Problem 6.59.
\subsection{Sources}
\begin{itemize}
\item UnsolvedMath local registry, permanent ID \(700616\), exact source problem code \texttt{catalogue entry}, titled ``Kourovka Notebook Problem 9.83''; supplied upstream revision \texttt{d371ee917ac9647fbb06cd6fc529deec6f3dda82}.
\item Attributed in the supplied record to V. P. Shunkov, proposed in 1984. The stated source is \emph{The Kourovka Notebook (Unsolved Problems in Group Theory)}, 9th issue (1984), Problem 9.83, as included in the 21st issue edition (2026), \texttt{arXiv:1401.0300v48}. The collection description names E. I. Khukhro and V. D. Mazurov as editors.
\item The statement refers to Kourovka Notebook Problem 6.59 for its technical hypothesis; that problem's text was not supplied.
\item The registry labels this problem ``open''. This is supplied source metadata, not an independently verified mathematical or literature assessment. The cited source has not been independently consulted for this draft.
\end{itemize}
\end{document}
